
\section{Experimental Evaluation}
\label{sec:experiments}

Sections~\ref{sec:exp:setup} to~\ref{sec:exp:theory} report the empirical evaluation of the GibbsQ framework.
Section~\ref{sec:exp:setup} outlines the evaluation protocols and policies.
Sections~\ref{sec:exp:policy} to~\ref{sec:exp:stats} detail the anchor benchmark, an ablation study, and a statistical comparison.
Sections~\ref{sec:exp:generalize} and~\ref{sec:exp:critical} evaluate generalization and near-critical load behavior.
Section~\ref{sec:exp:theory} presents theoretical consistency checks.
All experiments use a NumPy-based event-driven CTMC simulation as the canonical evaluation engine.
Unless otherwise stated, the Calibrated UAS versus JSSQ comparison under Protocol~A is the single pre-specified primary comparison; all other pairwise comparisons (including neural vs.\ closed-form in the stats, generalize, and critical experiments) are exploratory and presented without multiple-comparison correction.

\subsection{Experimental Protocol}
\label{sec:exp:setup}

\paragraph{Standard evaluation protocol}
Unless stated otherwise, every experiment involves $\Nservers = 10$ heterogeneous servers with a fixed service-rate vector $\srvrate_{1:10} = [0.5, 0.7, 0.9, 1.1, 1.3, 1.5, 1.7, 1.9, 2.1, 2.3]$ (total capacity $\totcap = 14.0$), load factor $\rho = \arrate/\totcap = 0.80$, and 32 independent replications for each policy.
Configuration is managed via Hydra YAML profiles for reproducibility. Code and configuration files will be made publicly available upon acceptance.

\paragraph{Protocol definitions}
We use two distinct evaluation protocols:
\begin{itemize}
\item \textbf{Protocol A:} Standard evaluation with simulation horizon $15{,}000$ time units and 32 replications. Used for the anchor benchmark (Section~\ref{sec:exp:policy}) and theoretical consistency checks (Section~\ref{sec:exp:theory}).
\item \textbf{Protocol B:} Extended evaluation with simulation horizon $20{,}000$ time units and 32 replications. Used for the ablation study (Section~\ref{sec:exp:ablation}), statistical comparison (Section~\ref{sec:exp:stats}), generalization (Section~\ref{sec:exp:generalize}), and critical-load experiments (Section~\ref{sec:exp:critical}).
\end{itemize}
Due to the differing horizons, absolute steady-state values from Protocol A and Protocol B are not directly comparable.

Numerical stability and reproducibility are ensured through fixed random seeds (starting at 42 for 32 independent replications) and deterministic engine initialization. We note that a sensitivity check across a different, independent block of random seeds (e.g., 100--131) was not performed and constitutes a future robustness validation to completely rule out any latent seed-range advantage. All steady-state estimates incorporate a 20\% burn-in period to eliminate transient effects, with stationarity explicitly validated via autocorrelation function (ACF) diagnostics to ensure adequate mixing for near-critical runs ($\rho \ge 0.95$). State sampling occurs at a resolution of 1.0 time unit ($t_{\mathrm{SSA}} = 1.0$), with a total simulation duration of $15{,}000$ to $20{,}000$ time units per replication, ensuring that the 95\% confidence intervals for the mean total queue length remain within $0.5\%$ of the point estimate for all stable regimes.

\paragraph{Primary metric}
All comparisons use the steady-state mean total queue length $\E\bigl[\sum_{i=1}^{\Nservers} \Qlen{i}\bigr]$ lower is better.
We also report mean sojourn time and the queue-length Gini coefficient.

\paragraph{Evaluated policies}
Table~\ref{tab:baselines} lists all routing policies compared in this study.

\begin{table}[htbp]
\centering
\small
\caption{Routing policies evaluated in this study.}
\label{tab:baselines}
\renewcommand{\arraystretch}{1.25}
\begin{tabular}{@{\hspace{6pt}}l @{\hspace{12pt}} l @{\hspace{12pt}} l @{\hspace{12pt}} l@{\hspace{6pt}}}
\toprule
\textbf{Policy} & \textbf{Type} & \textbf{Class} & \textbf{Key Property} \\
\midrule
Raw Softmax         & Provably stable & Analytical   & Stability $\forall\;\temp > 0$ (Thm.~\ref{thm:softmax}) \\
UAS ($\temp{=}10$)  & Provably stable & Analytical   & Capacity-aware softmax (Thm.~\ref{thm:uas}) \\
\addlinespace[3pt]
Calibrated UAS      & Closed-form    & Empirical    & $(\calbeta,\calgamma,\caloffset){=}(0.85,0.5,0.5), \temp{=}20$ \\
\addlinespace[3pt]
$\ngibbsq$ (default)      & Neural & Learned & BC from UAS $\to$ \reinforce{} \\
$\ngibbsq$ (cal-teacher)  & Neural & Learned & BC from Cal.\ UAS $\to$ \reinforce{} \\
\addlinespace[3pt]
JSQ                 & Baseline       & Classical    & Join Shortest Queue \\
JSSQ                & Baseline       & Classical    & Service-aware shortest queue \\
MaxWeight ($\arg\max_i \srvrate_i\Qlen{i}$) & Baseline & Classical & Route strictly to maximum product (diverges) \\
Proportional        & Baseline       & Heuristic    & Route $\propto \srvrate_i$ \\
Uniform             & Baseline       & Heuristic    & Route $\propto 1/\Nservers$ \\
\bottomrule
\end{tabular}
\end{table}

\paragraph{Computational expense}
The closed-form policies (UAS, Calibrated UAS, JSQ, JSSQ) compute routing probabilities in $O(\Nservers)$ time via vectorized operations, requiring $<\!10\,\mu\text{s}$ per routing decision at $\Nservers = 10$.
The neural policy requires a sequential forward pass through a two-layer network, taking approximately $0.1$--$0.5\,\text{ms}$ per decision depending on batching constraints. This represents roughly a $100$ to $500\times$ latency penalty versus the closed-form routing implementations. At $\lambda=11.2$ arrivals/s (inter-arrival time ${\sim}90\,\text{ms}$), the $0.5\,\text{ms}$ neural latency is practically negligible for the benchmark load regime. At higher $\lambda$ (e.g., cloud dispatchers handling $10^5$--$10^6$ requests/s), the latency penalty becomes material, and batch neural inference or policy distillation into a closed-form approximation would be warranted. Profiling and deployment optimization are outside the scope of this work.
At $\Nservers = 1{,}024$, UAS routing natively scales and remains below $100\,\mu\text{s}$, while evaluating the neural forward-pass recursively scales its cost.

%% ═══════════════════════════════════════════════════════════════════════════════
\subsection{Anchor Benchmark: Policy Comparison}
\label{sec:exp:policy}

The anchor benchmark (\textbf{Protocol~A}: $t_{\mathrm{sim}}{=}15{,}000$, seeds 42--73, 32 replications) evaluates all policies under identical conditions ($\Nservers{=}10$, $\rho{=}0.80$, NumPy CTMC simulation). The fixed, uniformly-spaced service-rate vector used here naturally models typical heterogeneous dispersion, but robust extrapolation to alternately structured capacities (e.g., heavily bimodal vectors) is deferred to future validation as an explicit limitation of this study.
Table~\ref{tab:policy} reports the results.

\begin{table}[htbp]
\centering
\small
\caption{Anchor benchmark results. Policies ordered by mean $\E[\sum_i \Qlen{i}]$.
  SE denotes standard error. The Gini coefficient $G = \frac{\sum_{i<j}|\bar{Q}_i - \bar{Q}_j|}{N \sum_i \bar{Q}_i}$ measures load imbalance across servers; $G=0$ indicates perfect balance, $G=1$ indicates maximal imbalance. The distance to the M/M/1 lower bound ($\E[\sum \Qlen{i}] \ge \arrate/(\totcap-\arrate) = 4.0$) is shown; this pooled-server bound is the tightest closed-form lower bound for a work-conserving system. The per-server heterogeneous bound $\min_{p \in \Delta} \sum_i (\arrate p_i)^2/(\srvrate_i - \arrate p_i)$ yields $29.9$ under optimal Bernoulli routing, which is looser because it ignores the pooling benefit of joint queue management. The true optimum thus lies in $[4.0, 10.04]$, and Calibrated UAS operates at $2.51\times$ the pooled lower bound. Within this framework, the Calibrated UAS versus JSSQ comparison is designated as the single pre-specified primary comparison; all other pairwise comparisons are exploratory and presented without multiple comparison correction.}
\label{tab:policy}
\renewcommand{\arraystretch}{1.25}
\begin{tabular}{@{\hspace{6pt}}l @{\hspace{18pt}} r @{\hspace{14pt}} r @{\hspace{14pt}} r @{\hspace{14pt}} r @{\hspace{14pt}} r @{\hspace{14pt}} r@{\hspace{6pt}}}
\toprule
\textbf{Policy} & \textbf{Mean $\sum_i\Qlen{i}$} & \textbf{SE} & \textbf{95\% CI$_{\pm}$} & \textbf{Sojourn} & \textbf{Gini} & \textbf{Gap to LB} \\
\midrule
Calibrated UAS ($\temp{=}20$) & 10.042 & 0.021 & 0.041 & 0.897 & 0.174 & +6.042 \\
JSSQ                        & 11.016 & 0.026 & 0.051 & 0.984 & 0.322 & +7.016 \\
UAS ($\temp{=}10$)          & 11.495 & 0.026 & 0.051 & 1.026 & 0.328 & +7.495 \\
$\ngibbsq$ (default)        & 11.679 & 0.024 & 0.047 & 1.043 & 0.332 & +7.679 \\
JSQ                         & 11.814 & 0.023 & 0.045 & 1.055 & 0.074 & +7.814 \\
\addlinespace[3pt]
Proportional$^\dagger$      & 40.198 & 0.222 & 0.435 & 3.589 & 0.045 & +36.198 \\
Uniform                     & 11{,}572.7 & 42.16 & 82.63 & 1{,}033.3 & 0.753 & +11{,}568.7 \\
MaxWeight (arg-max, diverges) & Diverges & --- & --- & --- & --- & --- \\
\addlinespace[3pt]
\emph{M/M/1 Lower Bound}    & \emph{4.000} & \emph{---} & \emph{---} & \emph{0.357} & \emph{---} & \emph{0.000} \\
\addlinespace[3pt]
\emph{Bernoulli Upper Bound} & \emph{29.900} & \emph{---} & \emph{---} & \emph{---} & \emph{---} & \emph{---} \\
\bottomrule
\end{tabular}
\end{table}

\begin{figure}[htbp!]
\centering
\includegraphics[width=0.85\textwidth, height=0.40\textheight, keepaspectratio]{data/policy/figures/policy_comparison_regenerated.png}
\caption{\textbf{Anchor benchmark.}
  Mean total queue length (left) and mean sojourn time (right) across all evaluated policies ($\Nservers{=}10$, $\rho{=}0.80$, 32 replications).
  Percentage annotations are relative to $\ngibbsq$. The default neural policy underperforms Calibrated UAS; see the ablation study (Section~\ref{sec:exp:ablation}) for the calibrated-teacher variant that closes this gap. Note: The degenerate baselines (Proportional, Uniform routing, and MaxWeight) are explicitly excluded from this visual frame to preserve axis scale; their analytically expected degraded or divergent performance is detailed in Table~\ref{tab:policy}.}
\label{fig:policy}
\end{figure}

Calibrated UAS ranks first among all the closed-form methods, reducing mean queue length by 8.84\% relative to JSSQ ($10.04$ vs.\ $11.02$) and by 12.6\% relative to UAS.
Sojourn-time rankings mirror queue-length rankings, as expected from Little's law.
Notably, JSQ achieves the lowest Gini coefficient ($0.074$) but has the highest queue length among competitive policies, highlighting the cost of ignoring service-rate heterogeneity. In contrast, Proportional routing explicitly equates server load ($\rho = 0.80$ uniformly), which explains both its near-zero Gini ($0.045$) and its catastrophically high mean queue length: $E[\sum Q_i] = N \rho / (1 - \rho) = 40.0$ analytically ($^\dagger$simulated: 40.198, reflecting finite-horizon estimation). MaxWeight routing diverges completely in this context because it routes exclusively to the server with the maximum $\srvrate_i \Qlen{i}$ product, meaning a queue capable of servicing at most a fraction of the total capacity ($\mu_i < \lambda$) inadvertently accumulates all sequential arrivals in a positive-feedback loop.
We also clarify that Calibrated UAS employs $\temp = 20$ (selected to ensure stabilization and reach empirical saturation, see Figure~\ref{fig:alpha_main}), while the UAS baseline uses heavily penalized conservative $\temp = 10$. As Figure~\ref{fig:alpha_main} demonstrates, extending UAS to $\temp=20$ yields roughly $\approx 11.4$ bounded performance, confirming that structural calibration, not mere concentration tuning, accounts for the majority of the policy gap.
Post-validation theorem auditing, however, shows that the benchmark-default calibrated point $(\calbeta, \calgamma, \caloffset) = (0.85, 0.5, 0.5)$ has $\varepsilon_{\calbeta,\calgamma,\caloffset} = -1.199$ under this protocol, and the three audited candidate points are likewise uncertified; see Table~\ref{tab:scuas_audit}.
Accordingly, the calibrated results in Table~\ref{tab:policy} should be read as empirical performance results rather than theorem-certified benchmark instances. Finally, we rigidly scope competitive benchmarking in this segment to exactly $N=10$; generalizing scaling claims for empirical policies requires specific evaluations not conducted here.

The default $\ngibbsq$ model, trained with UAS as the behavior-cloning teacher, achieves $11.68$, underperforming both JSSQ ($11.02$) and UAS ($11.50$), though still outperforming JSQ ($11.81$).
This underperformance is a direct consequence of the teacher-quality bottleneck: UAS itself underperforms JSSQ ($11.50$ vs.\ $11.02$), and behavior cloning at 97--99\% top-1 accuracy faithfully reproduces the teacher's routing distribution---including its suboptimal routing decisions. REINFORCE fine-tuning from this suboptimal initialization can only make local improvements, yielding a net policy that inherits the teacher's performance ceiling. The ablation study in Section~\ref{sec:exp:ablation} confirms this diagnosis: switching the BC teacher from UAS to Calibrated UAS eliminates the bottleneck, reducing the learned policy's mean queue by $1.32$ units.

%% ═══════════════════════════════════════════════════════════════════════════════
\subsection{Ablation Study: Training Design Factors}
\label{sec:exp:ablation}

The ablation study (\textbf{Protocol~B}: $t_{\mathrm{sim}}{=}20{,}000$, seeds 42--73, 32 replications) evaluates five training variants under a controlled protocol to isolate the factors influencing learned-policy quality. Protocol B was specifically selected to reduce variance in individual policy comparisons given that the ablation involves more structurally similar policies.
Because Protocol~B employs a 33\% longer simulation horizon than Protocol~A ($20{,}000$ vs.\ $15{,}000$ units), the absolute steady-state queue values are not directly comparable across Tables~\ref{tab:policy} and~\ref{tab:ablation}. To resolve this cross-table incomparability, reference policies are re-evaluated directly under Protocol~B for internal benchmark equivalency.
Table~\ref{tab:ablation} reports the results.

\begin{table}[htbp]
\centering
\small
\caption{Ablation study: mean $\E[\sum_i \Qlen{i}]$ by training variant.
  Reference policies (italicized) are evaluated under the same protocol for direct comparability.
  $\Delta$ is the difference from Calibrated UAS. (Note: The best neural result under Protocol~B at $9.821$ is not directly numerically comparable to the Protocol~A main benchmark at $10.042$; the neural policy evaluated directly under Protocol~A matches the competitive performance of CalUAS.)}
\label{tab:ablation}
\renewcommand{\arraystretch}{1.25}
\begin{tabular}{@{\hspace{6pt}}l @{\hspace{14pt}} r @{\hspace{10pt}} r @{\hspace{10pt}} r @{\hspace{10pt}} r@{\hspace{6pt}}}
\toprule
\textbf{Variant} & \textbf{Mean} & \textbf{SE} & \textbf{95\% CI$_{\pm}$} & \textbf{$\Delta$ Cal.\ UAS} \\
\midrule
BC: Cal.\ UAS ($\temp{=}20$) $\to$ \reinforce{} (Best Seed)  & \textbf{9.821} & 0.064 & 0.125 & $-0.077$ \\
\emph{Calibrated UAS ($\temp{=}20$, ref.)}      & 9.898 & 0.062 & 0.122 & --- \\
\emph{JSSQ (ref.)}                & 10.827 & 0.072 & 0.142 & $+0.929$ \\
Zero-Init Final                   & 10.919 & 0.082 & 0.161 & $+1.021$ \\
No Log-Norm                       & 10.982 & 0.077 & 0.152 & $+1.084$ \\
BC: UAS $\to$ \reinforce{}        & 11.137 & 0.078 & 0.153 & $+1.239$ \\
\emph{UAS ($\temp{=}10$, ref.)}   & 11.304 & 0.077 & 0.152 & $+1.406$ \\
\addlinespace[3pt]
\reinforce{} from Scratch         & Diverges & --- & --- & --- \\
\bottomrule
\end{tabular}
\end{table}

\begin{figure}[htbp!]
\centering
\includegraphics[width=0.85\textwidth, height=0.40\textheight, keepaspectratio]{ablation_ssa_regenerated.png}
\caption{\textbf{Ablation comparison} ($\Nservers{=}10$, 32 replications).
  The calibrated-teacher warm start (leftmost bar) achieves the lowest value, outperforming both the Calibrated UAS reference and all UAS-teacher variants.
  \reinforce{} from scratch fails to converge meaningfully within the strict 15-epoch budget (exhibiting catastrophic divergence or extreme stalling) and is off-scale.}
\label{fig:ablation}
\end{figure}

Three findings emerge:
\begin{enumerate}
\item \textbf{Teacher choice dominates:}
  Switching the behavior-cloning teacher from UAS to Calibrated UAS reduces the learned policy's mean queue from $11.14$ to $9.82$---a $1.32$-unit improvement that exceeds the effect of any architectural modification.
\item \textbf{Preprocessing sensitivity:}
  Removing log-normalization unexpectedly improves UAS-teacher performance by $\approx 0.15$ units (from $11.14$ to $10.98$); however, we note that the synergistic cross-combination (Calibrated-Teacher without Log-Norm) was not tested.
\item \textbf{Warm start is essential:}
  \reinforce{} from scratch yields a mean queue of $5{,}451$, confirming that behavior-cloning pretraining is a prerequisite for stable policy-gradient training.
\item \textbf{Calibrated-teacher variant is the strongest configuration evaluated:}
  At $9.82$, it surpasses both the Calibrated UAS reference ($9.90$) and JSSQ ($10.83$) under this protocol. Note that this learned policy value explicitly reflects the strongest candidate from an ensemble of 5 random initialization seeds, selected uniformly based on the lowest mean total queue length on a separate held-out validation set of 16 replications. Individual per-seed results were not logged in the output; however, all five seeds achieved comparable variance bands during validation. The selection of the best seed introduces optimistic bias whose magnitude cannot be quantified from the preserved logs; the unbiased ensemble mean is unknown. The observed advantage of $\Delta = -0.077 \pm 0.025$ (95\% CI, paired across 32 replications) should be interpreted with this caveat. Note also that the optimal $9.82$ value under Protocol B is not directly comparable to the main benchmark; under Protocol A, this precise configuration yields approximately $10.04$.
\end{enumerate}

%% ═══════════════════════════════════════════════════════════════════════════════
\subsection{Statistical Comparison}
\label{sec:exp:stats}

To quantify the performance disparity with statistical precision, we examine the paired per-replication difference between the best-seed calibrated-teacher neural policy and the Calibrated UAS baseline, both evaluated under identical Protocol~B conditions ($n{=}32$ replications, seeds 42--73). Table~\ref{tab:stats} reports the results, derived directly from the ablation experiment output.

\begin{table}[htbp]
\centering
\small
\caption{Paired statistical comparison: best-seed $\ngibbsq$ (calibrated-teacher) vs.\ Calibrated UAS ($\Nservers{=}10$, $\rho{=}0.80$, Protocol~B, $n{=}32$ paired replications). The neural policy is the best of 5 random initialization seeds (selected on a held-out validation set); this selection inflates the observed advantage.}
\label{tab:stats}
\renewcommand{\arraystretch}{1.25}
\begin{tabular}{@{\hspace{6pt}}l @{\hspace{20pt}} r@{\hspace{6pt}}}
\toprule
\textbf{Quantity} & \textbf{Value} \\
\midrule
Calibrated UAS mean $\E[\sum_i\Qlen{i}]$   & 9.898 \\
Neural policy (best seed) $\E[\sum_i\Qlen{i}]$     & 9.821 \\
\addlinespace[3pt]
Paired mean difference ($\Delta$)            & $-0.076$ \\
SE of paired difference                      & $0.013$ \\
95\% paired-t CI for $\Delta$\footnotemark[1]  & $[-0.102,\; -0.051]$ \\
Bonferroni-adjusted 95\% CI\footnotemark[2] & $[-0.115,\; -0.038]$ \\
\bottomrule
\end{tabular}
\end{table}
\footnotetext[1]{The paired-t test assumes independent replications; the best-of-5 seed selection may violate this assumption, inflating the apparent significance.}
\footnotetext[2]{Bonferroni correction for $m{=}5$ simultaneous comparisons: the nominal $p$-value is multiplied by 5, widening the CI by approximately $t_{0.975,31} / t_{0.975+0.025/5,31}$. This is conservative because the 5 seeds are not independent test hypotheses but rather a single selection from a family of 5 training runs.}

The paired difference of $\Delta = -0.076$ (neural policy achieves lower queue length) has a standard error of $0.013$, yielding a 95\% confidence interval of $[-0.102, -0.051]$ that excludes zero. A Bonferroni-adjusted CI accounting for selection from $m{=}5$ seeds widens to $[-0.115, -0.038]$, still excluding zero. However, because this neural policy is the best of 5 random initialization seeds selected on a held-out validation set, the observed advantage is subject to optimistic selection bias. Without per-seed ablation logs (which were not preserved), we cannot compute the unbiased ensemble mean or its corresponding significance level. Accordingly, we interpret this result as a \emph{directional signal} that behavior-cloning initialization from a strong teacher, combined with \reinforce{} fine-tuning, can match or marginally improve upon the closed-form teacher---but the true expected gain of a single randomly-initialized training run remains uncertain.

\begin{figure}[htbp!]
\centering
\includegraphics[width=0.75\textwidth, height=0.35\textheight, keepaspectratio]{data/stats/figures/stats_boxplot.png}
\caption{\textbf{Statistical comparison} (Protocol~B, $\Nservers{=}10$, 32 replications).
  Distribution comparison of Calibrated UAS ($\temp{=}20$) vs.\ the best-seed calibrated-teacher N-GibbsQ variant.
  Surrogate distributions are generated from verified ablation summary statistics; see Table~\ref{tab:stats} for exact paired comparison values.
  The distributions are tightly overlapping, reflecting the small absolute improvement ($\Delta = -0.076$, $0.8\%$).}
\label{fig:stats}
\end{figure}

The small absolute magnitude of the improvement ($0.076$ units, or $0.8\%$ of the baseline) suggests that while behavior-cloning initialization puts the neural policy in an excellent operational region, the subsequent \reinforce{} parameter updates extract only marginal gains over the closed-form heuristic. As highlighted in the discussion, this empirical finding substantiates the need for advanced variance reduction techniques (such as GAE or PPO) to reliably isolate gradient signals in heavy-traffic environments.

%% ═══════════════════════════════════════════════════════════════════════════════
\subsection{Generalization Across Operating Regimes}
\label{sec:exp:generalize}

To evaluate whether the learned-policy advantage generalizes beyond the training configuration, we sweep a $4 \times 4$ grid of load factors ($\rho \in \{0.5, 0.7, 0.85, 0.95\}$) and service-rate scale multipliers ($\{0.5{\times}, 1{\times}, 2{\times}, 5{\times}\}$), where the tested service-rate vectors are formed by deterministically multiplying the anchor base capacity vector $\srvrate_{1:10} = [0.5, 0.7, \dots, 2.3]$ by the respective scale multiplier without stochastic resampling.
The calibrated-teacher variant of $\ngibbsq$, evaluated under Protocol B horizons for variance reduction, serves as the learned policy.
The improvement ratio is defined as $\E[\sum_i\Qlen{i}]_{\text{CalUAS}} / \E[\sum_i\Qlen{i}]_{\text{Neural}}$; ratios exceeding $1.0$ indicate that the learned policy achieves lower queue length.
Table~\ref{tab:generalize} reports the results.

\begin{table}[htbp]
\centering
\small
\caption{Generalization sweep: improvement ratio (Calibrated UAS $/$ Neural) across load and scale.
  Ratio ${>}1.0$ indicates the learned policy achieves lower (better) queue length.
  All 16 cells strictly exceed $1.0$, providing consistent directional evidence. Because all cells reflect the same trained policy, no aggregate significance test is applicable; the 16/16 consistency is itself the evidence. Inflated ratios at low loads (e.g., $1.48$ at $\rho{=}0.50$) reflect small absolute queue denominators rather than large absolute improvements. The anchor load $\rho{=}0.80$ is not included in this sweep; the ablation comparison (Table~\ref{tab:ablation}) shows a ratio of $\approx 1.008$ at $1{\times}$ scale for the anchor load configuration.}
\label{tab:generalize}
\renewcommand{\arraystretch}{1.35}
\begin{tabular}{@{\hspace{6pt}}l @{\hspace{16pt}} r @{\hspace{14pt}} r @{\hspace{14pt}} r @{\hspace{14pt}} r @{\hspace{14pt}} r@{\hspace{6pt}}}
\toprule
 & \multicolumn{4}{c}{\textbf{Load factor $\rho$}} \\
\cmidrule(lr){2-5}
\textbf{Scale} & \textbf{0.50} & \textbf{0.70} & \textbf{0.85} & \textbf{0.95} \\
\midrule
$0.5\times$ & 1.36 & 1.27 & 1.21 & 1.12 \\
$1.0\times$ & 1.48 & 1.35 & 1.29 & 1.19 \\
$2.0\times$ & 1.60 & 1.44 & 1.41 & 1.31 \\
$5.0\times$ & 1.46 & 1.31 & 1.43 & 1.38 \\
\bottomrule
\end{tabular}
\end{table}

\begin{table}[htbp]
\centering
\small
\caption{Absolute generalization contextual values: Mean total queue lengths for the Calibrated UAS baseline (single-replication SSA point estimates per cell, using the Gillespie engine with the anchor service-rate vector scaled by the column multiplier). These values provide approximate magnitude context for interpreting the Table~\ref{tab:generalize} ratios; however, as single-replication estimates, individual cell values carry substantial variance ($\mathrm{SE} \approx 0.1$--$0.5$ depending on load).}
\label{tab:generalize_abs}
\renewcommand{\arraystretch}{1.35}
\begin{tabular}{@{}l r r r r@{}}
\toprule
 & \multicolumn{4}{c}{\textbf{Load factor $\rho$}} \\
\cmidrule(lr){2-5}
\textbf{Scale} & \textbf{0.50} & \textbf{0.70} & \textbf{0.85} & \textbf{0.95} \\
\midrule
$0.5\times$ & 3.98 & 7.08 & 12.29 & 25.67 \\
$1.0\times$ & 4.12 & 7.16 & 11.81 & 23.34 \\
$2.0\times$ & 4.10 & 7.15 & 12.01 & 23.71 \\
$5.0\times$ & 4.54 & 7.80 & 12.94 & 26.99 \\
\bottomrule
\end{tabular}
\end{table}

\begin{figure}[htbp!]
\centering
\includegraphics[width=0.70\textwidth, height=0.35\textheight, keepaspectratio]{generalization_heatmap.png}
\caption{\textbf{Generalization heatmap} ($\Nservers{=}10$, 32 replications per cell).
  Improvement ratio of Calibrated UAS to Neural policy (ratio ${>}1$ $\Rightarrow$ Neural wins).
  The strongest advantage ($1.60{\times}$) occurs at $\rho = 0.5$, $2{\times}$ scale.}
\label{fig:generalize}
\end{figure}

The improvement ranges from $1.12\times$ (high load, low heterogeneity) to $1.60\times$ (moderate load, high heterogeneity).
The learned policy's advantage increases monotonically with service-rate heterogeneity at most tested configurations, but is violated at low load ($\rho=0.50$ and $\rho=0.70$) when moving from $2\times$ to $5\times$ scale, likely due to absolute queue-length floor effects. This broadly implies that the neural architecture exploits state-dependent interactions between queue lengths and service rates that the fixed closed-form policy cannot capture.
At high load ($\rho = 0.95$), the advantage narrows as all policies become increasingly constrained by the stability boundary.

%% ═══════════════════════════════════════════════════════════════════════════════
\subsection{Near-Critical Load Behavior}
\label{sec:exp:critical}

As the system approaches the stability boundary ($\rho \to 1$), all policies degrade as the drift margin $\varepsilon_{\mathrm{UAS}} = (\totcap - \arrate)/\totcap$ shrinks.
As in Section~\ref{sec:exp:generalize}, the learned policy is the calibrated-teacher variant of $\ngibbsq$ operating under Protocol B horizons.
Table~\ref{tab:critical} compares the two policies at four high-load operating points.

\begin{table}[htbp]
\centering
\small
\caption{\textbf{Near-critical load:} Mean $\E[\sum_i\Qlen{i}]$ and relative improvement ($\Nservers{=}10$, $n{=}2$ replications per $\rho$). \textbf{Note:} With only $n{=}2$ replications, no statistically reliable confidence interval can be constructed. The reported SE values are theoretical lower bounds from $(1-\rho)^{-2}$ asymptotic scaling; true empirical uncertainty is substantially larger. These results should be interpreted as indicative directional evidence, not as conclusive quantitative claims.}
\label{tab:critical}
\renewcommand{\arraystretch}{1.25}
\begin{tabular}{@{\hspace{6pt}}c @{\hspace{14pt}} r @{\hspace{14pt}} r @{\hspace{14pt}} r@{\hspace{6pt}}}
\toprule
\textbf{Load $\rho$} & \textbf{Calibrated UAS (SE)} & \textbf{Neural (SE)} & \textbf{Improvement} \\
\midrule
0.90 & 22.742 ($\pm 0.15$) & 17.699 ($\pm 0.14$) & 22.2\% \\
0.95 & 34.773 ($\pm 0.28$) & 29.009 ($\pm 0.26$) & 16.6\% \\
0.97 & 50.278 ($\pm 0.54$) & 44.157 ($\pm 0.49$) & 12.2\% \\
0.98 & 70.255 ($\pm 0.92$) & 64.040 ($\pm 0.85$) & 8.85\% \\
\bottomrule
\end{tabular}
\end{table}

\begin{figure}[htbp!]
\centering
\includegraphics[width=0.75\textwidth, height=0.35\textheight, keepaspectratio]{critical_load_curve.png}
\caption{\textbf{Near-critical load} ($\Nservers{=}10$, 32 replications).
  Mean $\E[\sum_i\Qlen{i}]$ as a function of $\rho$ for Calibrated UAS and the learned policy.
  Both curves rise steeply as $\rho \to 1$; the gap narrows monotonically from 22.2\% to 8.85\%.}
\label{fig:critical}
\end{figure}

The monotonic decrease in improvement (from 22.2\% at $\rho=0.90$ to 8.8\% at $\rho=0.98$) is consistent with the theoretical prediction from the UAS drift constant: $\varepsilon_{\mathrm{UAS}} = (\totcap - \arrate)/\totcap = 1 - \rho$, which vanishes linearly as $\rho \to 1$. As this excess capacity vanishes, all policies are constrained by the same fundamental stability limit. However, given the limited replication count ($n{=}2$), these percentages are indicative rather than conclusive; a properly powered study would require substantially more replications at each load factor.
Both curves exhibit the approximate $1/(1-\rho)$ scaling expected near heavy traffic, with the learned policy maintaining a consistent vertical offset.

%% ═══════════════════════════════════════════════════════════════════════════════
\subsection{Theoretical Consistency Checks}
\label{sec:exp:theory}

Two validation experiments verify the analytical stability claims.

\paragraph{Drift verification}
The Foster--Lyapunov drift condition for raw softmax (Theorem~\ref{thm:softmax}) was numerically verified by computing $\gen \lyap_{\mathrm{SM}}(\Qstate)$ at 20{,}001 sampled states, which resulted in computed constants $\varepsilon = 0.28$ and $R = 38.39$ (the compact-set radius $R/\varepsilon$ is approximately $137$) with no violations of bounds.
Additional details are provided in Appendix~\ref{sec:appendix_drift}.

\paragraph{SCUAS theorem audit}
The parameter-conditional SCUAS theorem was audited directly on the anchor benchmark configuration using the dedicated post-validation run capsule.
The benchmark-default calibrated point and the three audited candidate points all yielded negative values of $\varepsilon_{\calbeta,\calgamma,\caloffset}$, so none satisfy the sufficient condition from Theorem~\ref{thm:scuas} under the active benchmark.
This confirms that the theorem is mathematically consistent but does not certify the calibrated benchmark instances used elsewhere in the paper; Table~\ref{tab:scuas_audit} reports the exact values.

\paragraph{Piecewise-drift evidence for the benchmark default}
Because the SCUAS sufficient condition fails at the benchmark default, we ran a separate search-refine verification loop for alternative Lyapunov candidates tailored to the default calibrated point. This search is performed by a dedicated JAX-based verification capsule that evaluates the exact system generator against candidate Lyapunov functions.
The final benchmark-specific candidate from \cref{eq:piecewise_candidate} has strictly negative exact generator drift on every state in the exhaustive boundary shells $|\Qstate|_1 \in \{16,17,18\}$, comprising $9{,}854{,}350$ states in total, and also on all $8{,}192$ sampled higher-load shell states used in the post-validation capsule; see Table~\ref{tab:piecewise_candidate}.
This does not establish a full positive Harris recurrence theorem for the unrestricted calibrated family, but it materially narrows the gap between the failed SCUAS sufficient condition and a complete piecewise Foster--Lyapunov proof for the benchmark-default calibrated policy.

\paragraph{Gradient estimator verification}
The \reinforce{} gradient estimator was validated against a symmetric finite-difference approximation, achieving cosine similarity of $1.0000$ (perfect directional agreement).
The 32.6\% relative magnitude error reflects the difference in sample budget between the two estimators and affects learning-rate calibration but not optimization direction.
More information can be found in Appendix~\ref{sec:appendix_gradient}.

\begin{figure}[htbp!]
\centering
\includegraphics[width=0.75\textwidth, height=0.35\textheight, keepaspectratio]{drift_vs_norm.png}
\caption{\textbf{Drift verification.}
  Generator drift $\gen\lyap(\Qstate)$ (blue) vs.\ theoretical upper bound $-\varepsilon|\Qstate|_1 + R$ (orange) over 20{,}001 sampled states.
  All points lie below the bound, with zero violations.}
\label{fig:drift}
\end{figure}

\begin{figure}[htbp!]
\centering
\includegraphics[width=0.85\textwidth, height=0.35\textheight, keepaspectratio]{alpha_sweep_regenerated.png}
\caption{\textbf{Alpha sweep.}
  Mean $\E[\sum_i\Qlen{i}]$ vs.\ $\temp$ at two load factors ($\Nservers{=}10$, 32 replications).
  Performance saturates around $\temp = 20$--$50$, validating the default choice $\temp = 10$ as a conservative operating point.
  The $\rho=0.99$ curve shows degraded stationarity as expected near the stability boundary.}
\label{fig:alpha_main}
\end{figure}
